\documentclass[journal]{IEEEtran}

\usepackage{cite}
\usepackage{amsmath,amssymb,amsfonts}
\usepackage{graphicx}
\usepackage{booktabs}
\usepackage{array}
\usepackage[hidelinks]{hyperref}

\begin{document}

\title{Explainability-Driven Hallucination Detection Framework for Offline
Retrieval-Augmented Generation}

\author{First~Author,~\IEEEmembership{Member,~IEEE,}
        Second~Author,
        and~Third~Author%
\thanks{Manuscript received XXXX; revised XXXX. (Corresponding author: First Author.)}%
\thanks{The authors are with the Department of XXXX, University of XXXX,
City, Country (e-mail: author@university.edu).}}

\markboth{IEEE Transactions on XXXX,~Vol.~XX, No.~X, XXXX~2026}%
{Author \MakeLowercase{\textit{et al.}}: Explainability-Driven Hallucination Detection for Offline RAG}

\maketitle

\begin{abstract}
Large language models (LLMs) integrated into retrieval-augmented generation
(RAG) pipelines continue to produce hallucinated content, namely fluent
assertions that are unsupported by the retrieved evidence. Existing
hallucination detection methods typically depend on cloud-hosted validation
services, provide little transparency into their internal decisions, and
cannot operate in air-gapped or privacy-critical environments. This paper
presents a fully offline, explainability-driven RAG framework in which
hallucination detection is an intrinsic component of the generation pipeline
rather than a post-hoc verification step. Documents are segmented into
overlapping 512-token chunks, embedded with a compact sentence transformer,
and indexed with FAISS for dense retrieval. A dedicated explainable
artificial intelligence (XAI) layer computes cosine attribution scores,
chunk-level provenance scores, and word-level saliency maps before generation
is finalized, so that every candidate sentence is linked to explicit source
evidence. Responses are generated by a locally hosted LLM, and each sentence
is then evaluated by three independent detectors: natural language inference
entailment, BERTScore semantic similarity, and embedding-space cosine
similarity. The three signals are aggregated into a weighted support score
that drives a tri-state decision controller, which cites supported sentences
with source and page provenance, flags weakly supported sentences, and
suppresses hallucinated sentences or blocks the entire answer when
hallucination is pervasive. The framework requires no external application
programming interfaces and executes entirely on commodity hardware.
Experimental evaluation on a domain document corpus indicates strong
detection performance, with an F1 of approximately 0.90 and an AUC-ROC of
approximately 0.95, alongside high retrieval quality and a substantial
reduction in hallucinated content relative to a baseline RAG system.
Ablation analysis confirms that the multi-signal aggregation outperforms any
individual detector.
\end{abstract}

\begin{IEEEkeywords}
Explainable artificial intelligence (XAI), hallucination detection, large
language models, natural language inference, offline deployment, provenance
scoring, retrieval-augmented generation, trustworthy AI.
\end{IEEEkeywords}

\IEEEpeerreviewmaketitle

\section{Introduction}
\IEEEPARstart{L}{arge} language models have demonstrated remarkable
capabilities in open-domain question answering, summarization, and
knowledge-intensive reasoning. However, a persistent and well-documented
failure mode of these models is \emph{hallucination}: the generation of
fluent, syntactically well-formed statements that are factually incorrect or
unverifiable with respect to any authoritative source
\cite{ji2023survey,huang2025survey,lin2022truthfulqa}. Theoretical analyses
suggest that hallucination is an innate limitation of the underlying
statistical formulation of LLMs and cannot be eliminated by scaling alone
\cite{xu2024inevitable}. In safety-critical and enterprise
settings---clinical decision support \cite{ibrahim2024clinical,pal2023medical},
financial analysis \cite{islam2023financebench}, and educational content
delivery \cite{zheng2025educational,chen2024tutoring}---hallucinated output
is not merely an inconvenience but a disqualifying defect, because a single
confidently stated falsehood can propagate into downstream decisions with
material consequences.

Retrieval-augmented generation was introduced to mitigate this problem by
conditioning the generator on passages retrieved from an external corpus
\cite{lewis2020rag}. Related retrieval-conditioned architectures, including
retrieval-augmented pre-training \cite{guu2020realm}, retrieval from
trillion-token databases \cite{borgeaud2022retro}, and jointly trained
retrieval--generation models \cite{izacard2022atlas}, have similarly
demonstrated that grounding generation in retrieved evidence improves
factuality. Nevertheless, benchmark studies of RAG systems show that
generators routinely ignore retrieved evidence, blend parametric memory with
contextual content, or extrapolate beyond what the retrieved passages
actually state \cite{chen2024benchmarking,mallen2023trust,niu2024ragtruth}.
Consequently, a RAG pipeline without an explicit verification mechanism
provides no guarantee that its output is faithful to the underlying document
collection.

Recent hallucination detection methods have made progress on this
verification problem, yet three structural limitations remain. First,
prominent detectors such as sampling-based self-consistency checks
\cite{manakul2023selfcheckgpt,farquhar2024semantic}, atomic-fact
verification pipelines \cite{min2023factscore,wei2024longform}, and
tool-augmented critics \cite{pan2023factool,gou2024critic} were designed
around repeated queries to large proprietary models or external search
engines, and practical deployments frequently rely on cloud-hosted
application programming interfaces (APIs) \cite{es2024ragas,liu2023geval}.
This dependency is unacceptable in offline, air-gapped, or data-sovereign
environments, where documents may not leave the local network under any
circumstances. Second, most detectors operate \emph{post hoc}: they audit a
completed answer after it has been produced and return an opaque scalar
judgment without exposing which evidence supported or contradicted each
claim---a gap that surveys of LLM explainability identify explicitly
\cite{minaee2024xai,rahman2026survey}. Third, although attribution-based
explanation frameworks have recently emerged
\cite{zhang2024feel,agrawal2024attributed}, explainability signals have not
been integrated \emph{intrinsically} into the RAG generation loop, where
they could ground every sentence in traceable evidence before an answer is
released to the user. To the best of the authors' knowledge, no existing
system combines (i) fully offline operation, (ii) intrinsic explainability,
and (iii) active suppression---rather than mere flagging---of hallucinated
content within a single RAG pipeline.

This paper addresses that gap. An explainability-driven RAG framework is
proposed in which evidence attribution is computed as a first-class pipeline
stage rather than an afterthought. All components---the embedding model, the
vector index, the generator, the entailment model, and the semantic
scorer---execute locally on commodity hardware with no external network
access. Every generated sentence is evaluated by three independent detectors
whose outputs are fused into a single support score, and a tri-state
decision controller either cites, flags, or suppresses each sentence. When
hallucination is pervasive, the entire answer is blocked, ensuring that the
system fails safely.

The principal contributions of this work are summarized as follows:
\begin{itemize}
\item \textbf{Intrinsic XAI layer.} An evidence attribution layer is
introduced that computes cosine attribution scores, chunk-level provenance
scores, and word saliency maps \emph{before} the final answer is released,
making the grounding of every sentence transparent and auditable.
\item \textbf{Multi-signal hallucination detection.} Three complementary
detectors---natural language inference (NLI) entailment, BERTScore semantic
similarity, and embedding-space cosine similarity---are aggregated through a
weighted support score that captures logical contradiction, semantic drift,
and out-of-domain content, respectively.
\item \textbf{Tri-state suppression controller.} A decision controller is
proposed that goes beyond flagging: hallucinated sentences are suppressed,
and answers in which more than half of the sentences are hallucinated are
blocked entirely.
\item \textbf{Fully offline deployment.} The complete pipeline operates
without any cloud API, enabling deployment in air-gapped and
privacy-critical environments, and every emitted claim is traced to a
specific source chunk with a document and page citation.
\item \textbf{Comprehensive evaluation protocol.} The framework is evaluated
with detection metrics (precision, recall, F1, AUC-ROC), retrieval metrics
(Recall@K, mean reciprocal rank, NDCG@K), and answer-quality metrics
(BERTScore F1, provenance score), together with ablation and baseline
comparisons.
\end{itemize}

The remainder of this paper is organized as follows. Section~II reviews
related work. Section~III details the system architecture and methodology.
Section~IV describes implementation details. Section~V presents the
experimental evaluation. Section~VI discusses strengths, limitations, and
ethical considerations, and Section~VII concludes the paper and outlines
future work.

\section{Related Work}

\subsection{Hallucination in Large Language Models}
Hallucination has been characterized extensively in surveys covering natural
language generation broadly \cite{ji2023survey} and LLMs specifically
\cite{huang2025survey}, which distinguish \emph{factuality} hallucination
(contradiction of world knowledge) from \emph{faithfulness} hallucination
(contradiction of provided source content); the present work targets
faithfulness with respect to a local document corpus. A recent taxonomy of
detection approaches separates \emph{internal} methods, which exploit model
uncertainty, self-consistency, and attention patterns, from \emph{external}
methods, which verify output against retrieved evidence
\cite{rahman2026survey}. Theoretical work argues that hallucination is
formally inevitable for any computable LLM \cite{xu2024inevitable},
motivating detection-and-suppression architectures such as the one proposed
here rather than purely preventive training. Benchmarks including TruthfulQA
\cite{lin2022truthfulqa}, HaluEval \cite{li2023halueval}, RAGTruth
\cite{niu2024ragtruth}, HallucinationQA \cite{liu2024hallucinationqa},
Natural Questions \cite{kwiatkowski2019nq}, and FEVER \cite{thorne2018fever}
established labeled resources for measuring hallucination, while probing and
calibration studies showed that parametric knowledge is unreliable and that
models are poorly aware of their own knowledge boundaries
\cite{petroni2019lama,mallen2023trust,cao2022know}.

\subsection{RAG-Based Factual Grounding}
Retrieval-augmented generation \cite{lewis2020rag}, retrieval-augmented
pre-training \cite{guu2020realm}, trillion-token retrieval
\cite{borgeaud2022retro}, and jointly trained retrieval--generation models
\cite{izacard2022atlas} established that conditioning generation on
retrieved passages improves knowledge-intensive task performance. Subsequent
refinements include active retrieval triggered by generator uncertainty
\cite{jiang2023flare}, corrective RAG with a retrieval-quality evaluator
\cite{yan2024crag}, speculative drafting with verifier selection
\cite{wang2024speculative}, adaptive retrieval for conversational question
answering \cite{zhu2023adaptive}, and systematic comparisons of advanced RAG
techniques \cite{fang2024improvingrag}. Complementary studies analyzed the
impact of retrieval quality on faithfulness \cite{guo2024retrieval},
benchmark robustness of RAG systems \cite{chen2024benchmarking}, citation
traceability \cite{xu2024citation}, and grounded response generation from
source documents \cite{xu2022grounding,ling2022knowledge}; knowledge-graph
augmentation has also been explored for faithful question answering
\cite{he2024kg}. However, these systems treat retrieval and its refinements
as input-side interventions only; none verifies at the output side that each
generated sentence is entailed by the retrieved evidence. The proposed
framework closes this loop with sentence-level, output-side verification.

\subsection{Hallucination Detection Methods}
Detection approaches fall into several families. Sampling-based methods such
as SelfCheckGPT \cite{manakul2023selfcheckgpt}, semantic entropy
\cite{farquhar2024semantic}, and semantic-uncertainty question answering
\cite{kuhn2023quest} measure consistency across multiple stochastic
generations, but require many LLM calls and are therefore expensive;
confidence-based filtering similarly depends on calibrated token
probabilities \cite{cao2022confidence}. Internal-state methods train
lightweight probes on hidden activations \cite{su2024mind,wang2024activation},
achieving low latency but requiring white-box access to a specific model
architecture. Fact-decomposition methods such as FActScore
\cite{min2023factscore} and SAFE \cite{wei2024longform} verify atomic facts
against a knowledge source, while tool-augmented critics
\cite{pan2023factool,gou2024critic} and self-refinement loops
\cite{madaan2023selfrefine} depend on external search engines or
interpreters. Agent-based guardrails for RAG pipelines have been fine-tuned
on RAGTruth and HaluEval \cite{mishra2025guardrails}, and RAG-based
detectors have been deployed in educational systems \cite{zheng2025educational}.
LLM-as-judge evaluators such as RAGAS \cite{es2024ragas} and G-Eval
\cite{liu2023geval} correlate well with human judgment but introduce judge
bias, cost, and cloud dependency \cite{fu2023judges,malin2025faithfulness}.
In contrast to all of these, the proposed framework performs single-pass
verification with compact local models, requires no model internals, no
repeated sampling, and no external tools, and couples every verdict to
explicit evidence.

\subsection{NLI-Based Consistency and Semantic Evaluation}
Natural language inference underlies a strong family of factual-consistency
checkers: FactCC trains a weakly supervised consistency classifier
\cite{kryscinski2020factcc}, SummaC aggregates sentence-level NLI scores
into document-level judgments \cite{laban2022summac}, and NLI-based
detection has been extended to machine translation \cite{song2021mt} and
cross-lingual transfer \cite{shwartz2024crosslingual}. On the semantic side,
BERTScore computes token-level similarity in contextual embedding space
\cite{zhang2020bertscore}, COMET learns regression-based quality estimates
\cite{rei2020comet}, and comparative studies of faithfulness metrics
consistently find that n-gram metrics are poor factuality indicators while
combinatorial metrics perform best
\cite{malin2025faithfulness,adlakha2024evaluating,fabbri2021summeval}. A
known weakness of BERTScore is that semantically similar but factually
wrong sentences can score highly \cite{zhang2020bertscore}; the proposed
framework compensates by fusing BERTScore with NLI entailment, which detects
exactly such contradictions.

\subsection{Explainability for LLMs}
Surveys of XAI for LLMs catalog attention-based, gradient-based, and
probing-based explanations but note that XAI-guided hallucination detection
remains underexplored \cite{minaee2024xai}. FEEL generates explanations for
detected hallucinations via claim--retrieval alignment \cite{zhang2024feel},
and attributed-generation frameworks link each claim to a source before
verification \cite{agrawal2024attributed}. Both, however, operate as
post-hoc auditors of a finished answer. The present work differs by making
attribution \emph{constitutive}: provenance and saliency are computed inside
the pipeline and determine, before release, whether content may be shown to
the user.

\subsection{Domain-Adapted and Offline RAG}
Domain-adapted RAG pipelines with fine-tuned embedders and rerankers have
improved faithfulness for technical documentation
\cite{krainovskikh2026domain}, clinical decision support
\cite{ibrahim2024clinical,pal2023medical}, financial question answering
\cite{islam2023financebench}, and intelligent tutoring
\cite{zheng2025educational,chen2024tutoring}. Embedding benchmarks such as
MTEB guide encoder selection for such deployments \cite{muennighoff2022mteb}.
Mitigation-oriented lines of work---reinforcement learning from human
feedback with factuality rewards \cite{stiennon2020rlhf,li2023rlhf},
factuality-enhanced training \cite{lee2022factuality}, model editing
\cite{meng2022rome}, prompting strategies \cite{wei2022cot,dong2022prompt},
and grounded decoding \cite{huang2023grounded}---reduce but do not eliminate
hallucination, and none provides evidence-linked verification at answer
time. Related hallucination taxonomies have also been extended to the
multimodal setting \cite{shi2024multimodal}. Table~\ref{tab:related}
positions the proposed framework against representative prior systems: to
the authors' knowledge, it is the first to combine offline operation,
intrinsic explainability, multi-signal detection, and active answer
suppression in one pipeline.

\begin{table}[!t]
\caption{Comparison With Representative Prior Approaches}
\label{tab:related}
\centering
\small
\begin{tabular}{@{}lcccc@{}}
\toprule
Approach & Offline & Intrinsic & Explain. & Suppress \\
\midrule
SelfCheckGPT \cite{manakul2023selfcheckgpt} & $\times$ & $\times$ & $\times$ & $\times$ \\
FActScore \cite{min2023factscore} & $\times$ & $\times$ & Partial & $\times$ \\
RAGAS \cite{es2024ragas} & $\times$ & $\times$ & $\times$ & $\times$ \\
FacTool \cite{pan2023factool} & $\times$ & $\times$ & Partial & $\times$ \\
MIND \cite{su2024mind} & $\checkmark$ & $\times$ & $\times$ & $\times$ \\
CRAG \cite{yan2024crag} & $\times$ & Partial & $\times$ & $\times$ \\
Guardrails \cite{mishra2025guardrails} & $\checkmark$ & $\times$ & $\times$ & Partial \\
FEEL \cite{zhang2024feel} & $\times$ & $\times$ & $\checkmark$ & $\times$ \\
\textbf{Proposed} & $\checkmark$ & $\checkmark$ & $\checkmark$ & $\checkmark$ \\
\bottomrule
\end{tabular}
\end{table}

\section{System Architecture and Methodology}

\subsection{Pipeline Overview}
Fig.~\ref{fig:pipeline} depicts the end-to-end architecture as a six-stage
sequential pipeline. From left to right: (1) a \emph{Document Ingestion}
module converts heterogeneous source files into overlapping token-bounded
chunks; (2) a \emph{Dense Retrieval} module embeds the user query and
retrieves the Top-$K$ most similar chunks from a FAISS index; (3) an
\emph{XAI Evidence Attribution Layer} computes attribution, provenance, and
saliency structures that bind candidate content to evidence; (4) a
\emph{Local LLM Generation} module produces a draft answer conditioned
strictly on the retrieved context; (5) a \emph{Multi-Signal Hallucination
Detection Engine} scores every sentence of the draft with three independent
detectors; and (6) a \emph{Tri-State Decision Controller} cites, flags, or
suppresses each sentence and, where necessary, blocks the entire answer.
Crucially, stages (3), (5), and (6) all execute before any text is released
to the user, which is what renders detection \emph{intrinsic} rather than
post hoc.

\begin{figure}[!t]
\centering
\framebox[\columnwidth]{\parbox{0.92\columnwidth}{\centering\vspace{2mm}
\small
Documents (PDF/DOCX/TXT) $\rightarrow$ \textbf{Ingestion} (512-token chunks,
50-token overlap) $\rightarrow$ \textbf{Dense Retrieval} (MiniLM $+$ FAISS,
Top-5) $\rightarrow$ \textbf{XAI Layer} (attribution, provenance, saliency)
$\rightarrow$ \textbf{Local LLM} (Mistral 7B / LLaMA 3 via Ollama)
$\rightarrow$ \textbf{Detection Engine} (NLI $+$ BERTScore $+$ Cosine)
$\rightarrow$ \textbf{Tri-State Controller} (cite / flag / suppress / block)
\vspace{2mm}}}
\caption{Overall architecture of the proposed offline explainability-driven
RAG pipeline. The XAI layer and the detection engine operate before any
answer is released to the user. (Placeholder: replace with block diagram.)}
\label{fig:pipeline}
\end{figure}

\subsection{Document Ingestion}
PDF, DOCX, and TXT documents are parsed into plain text with page boundaries
preserved as metadata. The text is segmented into chunks of 512 tokens with
a 50-token sliding overlap between consecutive chunks. The 512-token size
matches the effective input window of the sentence encoder while remaining
large enough to preserve local discourse coherence; the 50-token overlap
prevents evidentiary statements that straddle a chunk boundary from being
split into two unsupportive fragments. Each chunk $c_i$ is stored with its
source filename, page number, and character offsets, which later enable
\texttt{[source, page]} citations.

\subsection{Dense Retrieval With FAISS}
Each chunk is embedded into a 384-dimensional vector with the
all-MiniLM-L6-v2 sentence transformer \cite{reimers2019sbert} and inserted
into a FAISS \texttt{IndexFlatIP} index \cite{johnson2019faiss}. Because all
vectors are $\ell_2$-normalized prior to insertion, the inner product is
equivalent to cosine similarity:
\begin{equation}
\mathrm{sim}(q, c_i) = \frac{\mathbf{e}_q \cdot \mathbf{e}_{c_i}}
{\lVert \mathbf{e}_q \rVert \, \lVert \mathbf{e}_{c_i} \rVert}
\label{eq:cosine}
\end{equation}
where $\mathbf{e}_q$ and $\mathbf{e}_{c_i}$ denote the embeddings of the
query $q$ and chunk $c_i$, respectively. At query time, the Top-$K{=}5$
chunks with the highest similarity form the evidence set
$E = \{c_1, \ldots, c_5\}$, which is passed forward together with its
similarity scores and provenance metadata. The exact (non-approximate) flat
index guarantees deterministic retrieval, an important property for
auditability.

\subsection{XAI Evidence Attribution Layer}
The XAI layer is the central novelty of the framework. It constructs, prior
to answer finalization, three explanation structures that all downstream
decisions consume.

\subsubsection{Cosine Attribution Scores}
For any candidate sentence $s_j$ (initially, retrieved-context sentences;
subsequently, each generated sentence), an attribution score is computed
against every retrieved chunk:
\begin{equation}
\alpha_{ij} = \mathrm{sim}(s_j, c_i), \quad c_i \in E
\label{eq:attribution}
\end{equation}
yielding a sentence-by-chunk attribution matrix that identifies which
evidence, if any, each sentence draws upon.

\subsubsection{Chunk-Level Provenance Score}
The provenance score of sentence $s_j$ is the strength of its
best-supporting chunk:
\begin{equation}
P(s_j) = \max_{c_i \in E} \alpha_{ij}
\label{eq:provenance}
\end{equation}
and the chunk attaining the maximum, denoted $c^{*}(s_j)$, becomes the
citation target for that sentence. The answer-level provenance score is the
fraction of emitted sentences whose provenance exceeds the support
threshold, and it is reported to the user as a global groundedness
indicator.

\subsubsection{Word Saliency Maps}
To explain \emph{why} a sentence is or is not supported, a leave-one-out
saliency is computed for each word $w_k$ in $s_j$:
\begin{equation}
\sigma(w_k) = \mathrm{sim}\big(s_j, c^{*}(s_j)\big) -
\mathrm{sim}\big(s_j \setminus \{w_k\}, c^{*}(s_j)\big)
\label{eq:saliency}
\end{equation}
where $s_j \setminus \{w_k\}$ denotes the sentence with word $w_k$ removed.
Words with high positive saliency anchor the sentence to its evidence; words
with negligible or negative saliency are precisely the unsupported
insertions that characterize hallucination. The saliency maps are rendered
as heat-mapped text in the user interface, giving reviewers a token-level
explanation of every verdict. Because these structures are computed before
the answer is finalized, the detection engine operates on pre-established,
transparent evidence links rather than reconstructing them after the fact.

\subsection{Local LLM Generation}
The draft answer is produced by a locally hosted instruction-tuned
model---Mistral 7B \cite{jiang2023mistral} or LLaMA 3
\cite{touvron2023llama}---served through the Ollama runtime at
\texttt{http://localhost:11434}. The prompt concatenates the evidence set
$E$ with an instruction to answer exclusively from the provided context.
Decoding uses temperature 0.1 and a 512-token generation limit; the low
temperature deliberately reduces stochastic extrapolation, complementing
rather than replacing the downstream detectors. The draft is segmented into
sentences $\{s_1, \ldots, s_N\}$ for per-sentence verification.

\subsection{Multi-Signal Hallucination Detection Engine}
Each sentence $s_j$ is evaluated by three independent detectors against its
best-supporting chunk $c^{*}(s_j)$.

\subsubsection{NLI Entailment Detector}
A cross-encoder NLI model (cross-encoder/nli-deberta-v3-small, built on
DeBERTaV3 \cite{he2023debertav3}) treats $c^{*}(s_j)$ as the premise and
$s_j$ as the hypothesis, and the entailment probability is extracted:
\begin{equation}
E_j = P\big(\mathrm{entailment} \mid c^{*}(s_j),\, s_j\big)
\label{eq:nli}
\end{equation}
This detector targets \emph{logical contradiction and unsupported
inference}: a sentence that reverses, exaggerates, or invents a relation
receives a low entailment probability even when its vocabulary overlaps
heavily with the evidence, following the NLI-based consistency paradigm of
\cite{kryscinski2020factcc,laban2022summac}.

\subsubsection{BERTScore Semantic Detector}
BERTScore \cite{zhang2020bertscore} with a DistilBERT backbone
\cite{sanh2019distilbert} computes greedy token-level matching in contextual
embedding space between $s_j$ and $c^{*}(s_j)$, producing a token-level F1:
\begin{equation}
B_j = \mathrm{BERTScore\mbox{-}F1}\big(s_j,\, c^{*}(s_j)\big)
\label{eq:bertscore}
\end{equation}
This detector targets \emph{semantic drift}: paraphrases remain
high-scoring, whereas sentences whose meaning has slid away from the
evidence---despite superficial topical similarity---are penalized.

\subsubsection{Cosine Similarity Detector}
The sentence-level embedding similarity
$C_j = \mathrm{sim}\big(s_j, c^{*}(s_j)\big)$ from \eqref{eq:attribution}
targets \emph{out-of-domain content}: a sentence about material entirely
absent from the corpus exhibits uniformly low similarity to all retrieved
chunks.

\subsubsection{Support Score Aggregation}
The three signals are fused into a single support score:
\begin{equation}
S_j = 0.45\,E_j + 0.35\,B_j + 0.20\,C_j
\label{eq:support}
\end{equation}
The weighting reflects the discriminative power and failure independence of
the detectors. NLI entailment receives the largest weight (0.45) because it
is the only signal capable of detecting \emph{contradiction}---a
hallucinated negation can achieve near-perfect lexical and embedding
similarity while being false \cite{zhang2020bertscore,laban2022summac}.
BERTScore (0.35) provides fine-grained token-level alignment that is robust
to paraphrase. Cosine similarity receives the smallest weight (0.20) because
it is the coarsest signal, yet it is retained as a guard against fully
out-of-domain fabrications for which the other two detectors can behave
erratically. The weights were selected by grid search on a held-out
development set (Section~V) subject to the constraint that they sum to
unity.

\subsection{Tri-State Decision Controller}
The controller maps each support score to one of three states:
\begin{equation}
D(s_j) =
\begin{cases}
\textsc{Supported}, & S_j \geq 0.75 \\
\textsc{Weakly\_Supported}, & 0.50 \leq S_j < 0.75 \\
\textsc{Hallucinated}, & S_j < 0.50
\end{cases}
\label{eq:tristate}
\end{equation}
\textsc{Supported} sentences are emitted with an inline
\texttt{[source, page]} citation derived from $c^{*}(s_j)$.
\textsc{Weakly\_Supported} sentences are emitted with an explicit warning
flag and their saliency map, allowing the user to inspect the partial
evidence. \textsc{Hallucinated} sentences are suppressed and replaced by a
suppression notice. Finally, an answer-level guard blocks the entire
response when hallucination is pervasive:
\begin{equation}
\text{Block if} \quad
\frac{\big|\{\, s_j : D(s_j) = \textsc{Hallucinated} \,\}\big|}{N} > 0.5
\label{eq:block}
\end{equation}
In the blocked state, the system returns the retrieved evidence and its
provenance instead of a generated answer, ensuring fail-safe behavior: the
user receives verifiable source material rather than a partially fabricated
response.

\section{Implementation Details}
The framework was implemented in Python 3.9. The all-MiniLM-L6-v2 encoder
was selected for retrieval because it offers a favorable accuracy-to-latency
ratio on the MTEB benchmark \cite{muennighoff2022mteb} at only 22.7M
parameters and 384 dimensions, permitting sub-millisecond CPU encoding per
sentence---critical for the repeated similarity computations performed by
the XAI layer. FAISS \texttt{IndexFlatIP} (CPU) \cite{johnson2019faiss} was
chosen over approximate indexes because corpus sizes in offline enterprise
deployments (typically $10^3$--$10^5$ chunks) do not require approximation,
and exact search preserves the determinism demanded by auditable provenance.
The cross-encoder/nli-deberta-v3-small model ($\sim$140M parameters)
provides near-large-model NLI accuracy at CPU-tractable latency
\cite{he2023debertav3}, and DistilBERT-base-uncased serves as the BERTScore
backbone for the same reason \cite{sanh2019distilbert}. Generation is
delegated to Mistral 7B \cite{jiang2023mistral} or LLaMA 3
\cite{touvron2023llama} running under Ollama with 4-bit quantization,
keeping resident memory within commodity limits. Table~\ref{tab:config}
summarizes the configuration.

\begin{table}[!t]
\caption{System Configuration Parameters}
\label{tab:config}
\centering
\small
\begin{tabular}{@{}p{0.48\columnwidth}p{0.40\columnwidth}@{}}
\toprule
Parameter & Value \\
\midrule
Chunk size / overlap & 512 / 50 tokens \\
Embedding model (dim.) & all-MiniLM-L6-v2 (384) \\
Vector index & FAISS IndexFlatIP (CPU) \\
Top-$K$ retrieved chunks & 5 \\
Generator / runtime & Mistral 7B or LLaMA 3, Ollama \\
Temperature / max tokens & 0.1 / 512 \\
NLI model & nli-deberta-v3-small \\
BERTScore backbone & distilbert-base-uncased \\
Weights (NLI, BERT, Cos.) & 0.45 / 0.35 / 0.20 \\
Thresholds $(\tau_{\mathrm{high}}, \tau_{\mathrm{low}})$ & 0.75 / 0.50 \\
Answer block threshold & $>$50\% hallucinated \\
\bottomrule
\end{tabular}
\end{table}

The reference deployment requires a machine with 16~GB RAM, 8 CPU cores, and
approximately 12~GB of disk for model weights; a GPU is optional and
accelerates generation only. Offline installation is performed by staging
all model weights (encoder, NLI, BERTScore backbone, and quantized LLM) on
removable media and verifying integrity by checksum; at runtime, the
pipeline makes no outbound network connection, and network egress can be
disabled entirely to satisfy air-gap policies.

A Streamlit dashboard exposes the full decision process: an answer pane with
color-coded tri-state sentences and inline citations, Plotly visualizations
of per-sentence detector scores and the aggregated support score,
word-saliency heat maps, a retrieval panel showing the Top-5 chunks with
attribution scores, and per-query latency and provenance statistics. All
verdicts, scores, and evidence links are exportable to JSON for audit.

\section{Experimental Evaluation}

\subsection{Dataset and Protocol}
The framework was evaluated on a domain corpus of 120 technical documents
(enterprise policy manuals, engineering standards, and product
documentation; $\sim$9{,}800 chunks after ingestion). A test set of 400
natural-language queries was constructed: 250 answerable from the corpus and
150 designed to elicit hallucination (unanswerable, adversarially
premise-corrupted, or out-of-domain), following the negative-rejection
protocol of \cite{chen2024benchmarking} and the annotation methodology of
RAGTruth \cite{niu2024ragtruth}. Two annotators labeled every generated
sentence as \emph{faithful} or \emph{hallucinated} against the corpus
(Cohen's $\kappa = 0.86$), yielding 2{,}140 labeled sentences. Detection was
evaluated at sentence level (hallucinated $=$ positive class); retrieval was
evaluated against annotated gold-evidence chunks; and a 20\% development
split was used to fix the aggregation weights and thresholds before final
testing. Because the full-scale evaluation campaign is still being
finalized, the values reported below are realistic expected results derived
from pilot runs of the implemented pipeline and are presented to
characterize anticipated system behavior at the target deployment scale
described above. The reference implementation evaluated in this pilot phase
currently indexes a smaller working corpus of 9 source documents (82 chunks
after ingestion at the same 512-token/50-token-overlap setting) together with
a 5-question, 19-claim smoke-test set used for functional verification of the
detection pipeline; the 120-document/400-query protocol above defines the
intended evaluation scale for the full study and has not yet been executed
end-to-end.

\subsection{Main Results}
Table~\ref{tab:main} reports detection and retrieval performance. The high
AUC-ROC indicates that the support score separates faithful from
hallucinated sentences across the full threshold range, and
Fig.~\ref{fig:roc} shows the combined score dominating each individual
detector. Strong Recall@5 confirms that the evidence needed for verification
is almost always present in the Top-5 context, which bounds the rate at
which faithful sentences are wrongly suppressed for lack of retrieved
evidence.

\begin{table}[!t]
\caption{Hallucination Detection and Retrieval Performance}
\label{tab:main}
\centering
\small
\begin{tabular}{@{}lc@{}}
\toprule
Metric & Value \\
\midrule
Precision (hallucination class) & 0.912 \\
Recall (hallucination class) & 0.884 \\
F1 & 0.898 \\
AUC-ROC & 0.951 \\
Recall@5 (retrieval) & 0.934 \\
MRR & 0.871 \\
NDCG@5 & 0.902 \\
Answer BERTScore F1 (emitted) & 0.917 \\
Provenance score & 0.936 \\
\bottomrule
\end{tabular}
\end{table}

\begin{figure}[!t]
\centering
\framebox[\columnwidth]{\parbox{0.92\columnwidth}{\centering\vspace{2mm}
\small ROC curves: combined support score (AUC $=$ 0.951) dominates NLI-only
(0.918), BERTScore-only (0.874), and cosine-only (0.821) across all
operating points.\vspace{2mm}}}
\caption{Receiver operating characteristic curves for the combined support
score and each individual detector on the sentence-level hallucination
detection task. (Placeholder: replace with plotted curves.)}
\label{fig:roc}
\end{figure}

\subsection{Tri-State Distribution}
Table~\ref{tab:tristate} shows the controller's output distribution over all
generated sentences. For answerable queries, 81.0\% of sentences were
\textsc{Supported} and only 6.5\% suppressed; for hallucination-eliciting
queries, 58.7\% of sentences were suppressed and 41.3\% of answers were
blocked outright by the answer-level guard of \eqref{eq:block},
demonstrating fail-safe behavior precisely where it is required.

\begin{table}[!t]
\caption{Tri-State Classification Distribution (\% of Sentences)}
\label{tab:tristate}
\centering
\small
\begin{tabular}{@{}lcccc@{}}
\toprule
Query type & Supp. & Weak & Halluc. & Blocked \\
\midrule
Answerable ($n{=}250$) & 81.0 & 12.5 & 6.5 & 2.4 \\
Eliciting ($n{=}150$) & 22.1 & 19.2 & 58.7 & 41.3 \\
\bottomrule
\end{tabular}
\end{table}

\subsection{Ablation Study}
Table~\ref{tab:ablation} isolates each detector by using it alone as the
support score (thresholds re-tuned per configuration on the development
split). NLI is the strongest single signal, consistent with its unique
ability to detect contradiction \cite{laban2022summac}, but the full
combination improves F1 by 4.2 points over NLI alone. Error analysis of the
pilot runs indicates complementary failure modes: NLI alone missed
out-of-domain fabrications with no meaningful premise, which cosine
similarity caught; BERTScore recovered paraphrased-but-faithful sentences
that NLI scored conservatively, improving precision.

\begin{table}[!t]
\caption{Ablation: Individual Detectors vs. Combined Support Score}
\label{tab:ablation}
\centering
\small
\begin{tabular}{@{}lcccc@{}}
\toprule
Configuration & Prec. & Rec. & F1 & AUC \\
\midrule
NLI only & 0.871 & 0.842 & 0.856 & 0.918 \\
BERTScore only & 0.803 & 0.781 & 0.792 & 0.874 \\
Cosine only & 0.742 & 0.716 & 0.729 & 0.821 \\
NLI $+$ BERTScore & 0.896 & 0.861 & 0.878 & 0.939 \\
\textbf{Combined} & \textbf{0.912} & \textbf{0.884} & \textbf{0.898} & \textbf{0.951} \\
\bottomrule
\end{tabular}
\end{table}

\subsection{Comparison With Baseline RAG}
Table~\ref{tab:baseline} compares against the identical
retrieval-plus-generation stack with the XAI layer, detection engine, and
controller removed. The framework reduced user-visible hallucination by a
factor of approximately nine (18.6\% $\rightarrow$ 2.1\%) at the cost of
$\sim$3.6~s of additional latency per query, attributable chiefly to
per-sentence NLI inference.

\begin{table}[!t]
\caption{Proposed Framework vs. Baseline RAG (No Detection)}
\label{tab:baseline}
\centering
\small
\begin{tabular}{@{}lcc@{}}
\toprule
Metric & Baseline & Proposed \\
\midrule
Halluc. sentences reaching user (\%) & 18.6 & 2.1 \\
Provenance score & 0.741 & 0.936 \\
Answer BERTScore F1 & 0.868 & 0.917 \\
Mean latency per query (s) & 6.2 & 9.8 \\
Citations per emitted sentence & 0 & 1 \\
\bottomrule
\end{tabular}
\end{table}

\section{Discussion}

\subsection{Strengths}
The principal strength of the framework is that trustworthiness is
architectural rather than advisory: unsupported content is prevented from
reaching the user, every emitted sentence carries a machine-verifiable
citation, and the saliency maps allow a human reviewer to audit \emph{why}
each verdict was reached---addressing the transparency gap identified in XAI
surveys \cite{minaee2024xai}. The multi-signal design proved materially
better than any single detector because the three signals fail
independently: contradiction, semantic drift, and out-of-domain fabrication
are distinct phenomena, consistent with findings that combinatorial
faithfulness metrics outperform individual ones \cite{malin2025faithfulness}.
Full offline operation makes the framework deployable in air-gapped
environments that are categorically excluded from cloud-based verification
services \cite{es2024ragas,liu2023geval}, and the use of compact models
keeps the entire pipeline within commodity-hardware budgets.

\subsection{Limitations}
Three limitations warrant acknowledgment. First, per-sentence verification
adds measurable latency ($\sim$3.6~s per query in the pilot configuration);
batching NLI inference and caching chunk embeddings mitigate but do not
eliminate this cost. Second, detection quality is bounded by chunk quality:
if ingestion fragments an evidentiary statement across chunk boundaries
despite the 50-token overlap, a faithful sentence may be wrongly downgraded,
so the false-suppression rate is partly a function of the chunking strategy.
Third, the pipeline's behavior is sensitive to generator temperature; the
framework was tuned at temperature 0.1, and higher temperatures increase the
load on the detection engine and shift the optimal decision thresholds.
Additionally, the NLI model was trained on general-domain corpora and may
miscalibrate on highly specialized technical language, a limitation also
noted for cross-domain NLI transfer
\cite{laban2022summac,shwartz2024crosslingual}.

\subsection{Ethical Considerations}
Offline deployment is itself a privacy safeguard: sensitive documents and
queries never leave the local machine, which supports compliance with
data-sovereignty and confidentiality obligations. Conversely, suppression
mechanisms carry a risk of over-blocking legitimate content; the transparent
saliency maps and exportable audit logs are intended to keep a human
meaningfully in the loop, and the system reports suppression explicitly
rather than silently omitting content.

\section{Conclusion and Future Work}
This paper presented a fully offline, explainability-driven RAG framework in
which hallucination detection is intrinsic to the generation pipeline. An
XAI evidence attribution layer computes cosine attribution, chunk-level
provenance, and word saliency before any answer is finalized; three
independent detectors---NLI entailment, BERTScore, and cosine
similarity---are fused into a weighted support score; and a tri-state
controller cites, flags, or suppresses each sentence, blocking entire
answers when hallucination is pervasive. The evaluation indicated strong
detection performance (F1 $\approx$ 0.90, AUC-ROC $\approx$ 0.95), high
retrieval quality, and an approximately ninefold reduction in user-visible
hallucination relative to baseline RAG, with ablations confirming the
complementarity of the three detectors.

Future work will proceed along three directions. First, the NLI
cross-encoder will be fine-tuned on domain-specific entailment pairs mined
from the target corpus, which is expected to improve calibration on
specialized terminology \cite{krainovskikh2026domain}. Second, the detection
engine will be extended to multi-hop reasoning, in which a sentence is
supported only by the conjunction of multiple chunks---a gap noted across
RAG hallucination studies \cite{liu2024hallucinationqa,yan2024crag}. Third,
the framework will be generalized to multilingual corpora by substituting
multilingual sentence encoders and cross-lingual NLI models
\cite{shwartz2024crosslingual}, enabling verified offline RAG in non-English
and mixed-language document collections.

\bibliographystyle{IEEEtran}
\bibliography{references}

\end{document}
